\chapter{Onychomycosis}

\medskip
\noindent\\textit{\\textbf{Under review.} This chapter is an AI-assisted educational draft; it is not a substitute for professional medical advice.}

\section*{Definition and Overview}
a fungal infection of a fingernail or toenail that causes thickening, discoloration, brittleness, or separation from the nail bed. Individual assessment is important because symptoms and treatment vary with the cause.

\section*{Clinical Features}
Yellow, white, or brown discoloration, thick or crumbly nails, subungual debris, and discomfort in footwear.

\section*{When to Seek Medical Care}
Seek medical review for new, persistent, worsening, or function-limiting symptoms. Seek emergency help for severe breathing difficulty, collapse, sudden neurological change, severe chest or abdominal pain, or any rapidly worsening illness.

\section*{Diagnosis and Investigations}
Nail examination and, before prolonged oral treatment, confirmation with microscopy, culture, or other laboratory testing because trauma and psoriasis can look similar.

\section*{Management}
Keep feet dry, treat coexisting tinea pedis, and use topical or oral antifungal treatment selected by a clinician; improvement takes months as the nail grows.

\section*{Complications}
Possible complications depend on the cause and may include progression of the condition, effects on other organs, treatment adverse effects, and reduced quality of life. Early assessment can reduce preventable harm.

\section*{Prognosis and Outcomes}
Outcome depends on the underlying cause, severity, complications, and response to treatment. Discuss individual prognosis with the treating clinician.

\section*{Prevention and Self-Care}
Follow the treatment plan, keep relevant appointments, avoid known triggers, protect affected areas from injury, and seek help early if symptoms worsen. Do not start or stop prescription treatment without clinical advice.

\clearpage
